\subsection{Arithmetic folding with the original quotient norm}
\label{global-gram-continuation}

Retain an existing arithmetic setting $X$ and an original word $u$.
Write $E/F$ for its CM extension, $n$ for its moving dimension,
$G=U(V)$, $\Gamma=G(F)$ diagonally embedded in $G(\mathbb A_F)$,
and $S_u=N_u\rtimes H_u$.  The group $H_u$ is the unitary group of
the original positive definite orthogonal complement.  It is compact
at every infinite place.  All Haar measures and the quotient measure
$dq_u$ of Theorem~\ref{actual-seed-decay-theorem} remain fixed.
Put
\[
 M(g)=\prod_{v\mid\infty}\|g_v\|_{\rm op},\qquad
 \nu_u(g)=\inf_{s\in S_u(F_\infty)}M(sg),\qquad
 N_X=4n^2\sqrt{2\max_vq_v}.
\]
For compact $\Omega\subset G(\mathbb A_F)$ define
\[
 m_\Omega(x)=\sup_{y\in G(\mathbb A_F)}
 \#\{\gamma\in\Gamma:x^{-1}\gamma y\in\Omega\}.
\]
This is the original return multiplicity.  Constants below may depend
on the entire fixed arithmetic setting.  They are not uniform estimates
for the final growing family.\cite{global-gram-continuation-source}

\begin{theorem}[Arithmetic return multiplicity]
\label{global-gram-multiplicity}
For every pair of compact sets $\Omega,C\subset G(\mathbb A_F)$
there is a finite $A_{X,\Omega,C}\ge0$ such that
\begin{equation}\label{global-gram-mult-bound}
 m_\Omega(c(g,1))\le A_{X,\Omega,C}M(g)^{4n^2}
 \qquad(c\in C,\ g\in G(F_\infty)).
\end{equation}
Consequently the original condition \texttt{MultPoly}, with exponent
$N_X$, holds.  Neither $[F:\mathbb Q]$ nor $n$ is assumed bounded.
\end{theorem}

\begin{proof}
Fix $x=c(g,1)$ and $y$.  If the return set is nonempty, choose
$\gamma_0$ in it.  The injective map $\gamma\mapsto\gamma\gamma_0^{-1}$
sends that set into
\[
 \Gamma\cap x\Omega\Omega^{-1}x^{-1}.
\]
The finite components of the latter compact set lie in the fixed
compact set $C_f\Omega_f\Omega_f^{-1}C_f^{-1}$.  In a fixed
$E$-basis of $V$, its matrix entries therefore belong to
$a^{-1}\widehat{\mathcal O}_E$ for one positive integer $a$.
A rational matrix in this set has entries in $a^{-1}\mathcal O_E$.
Thus all candidates belong to the single Minkowski lattice
$\Lambda=(a^{-1}\mathcal O_E)^{n^2}$ in
$\prod_{v\mid\infty}\operatorname{Mat}_n(\mathbb C)$.

In the original normalized real coordinates, $g_v^{-1}=J_vg_v^*J_v$
with $J_v$ a diagonal matrix of signs.  Hence
$\|g_v^{-1}\|_{\rm op}=\|g_v\|_{\rm op}$ and
$\|g_v\|_{\rm op}\ge1$.  Compactness of $C$ and $\Omega$, including
the fixed change from the $E$-basis to these real coordinates, gives
constants $B_v$ for which every candidate matrix has Frobenius norm
at most $B_v\|g_v\|_{\rm op}^2$ in the $v$-block.

Choose a bounded half-open fundamental parallelepiped $P$ of
$\Lambda$, with $v$-block radius $p_v$.  Its translates by distinct
lattice points have disjoint interiors.  The translates associated
with our candidates lie in the product of balls of radii
$(B_v+p_v)\|g_v\|_{\rm op}^2$.  Comparing Euclidean volumes gives
\[
 \#\bigl(\Gamma\cap x\Omega\Omega^{-1}x^{-1}\bigr)
 \le \frac{\prod_v\operatorname{vol}_{2n^2}(B(0,B_v+p_v))}
                {\operatorname{covol}(\Lambda)}
       \prod_v\|g_v\|_{\rm op}^{4n^2}.
\]
This is independent of $y$.  It proves \eqref{global-gram-mult-bound};
empty compact sets cause no difficulty.  If some $q_v>0$, then
$N_X\ge4n^2$.  If every $q_v=0$, every $g_v$ is unitary and
$M(g)=1$, so the required exponent $N_X=0$ also works.
The multiplicity is Borel measurable.  For each integer $j$, the
condition $m_\Omega(x)\ge j$ is a countable union over distinct
$\gamma_1,\ldots,\gamma_j\in\Gamma$ of closed conditions: eliminate
$y$ using the first return, and project the resulting closed set
in $G(\mathbb A_F)\times\Omega$ along its compact second factor.
\end{proof}

Here is its exact application to the original scalar seeds.  Write
$K=K_2$, $x_u\in G(\mathbb A_{F,f})$ for the prescribed finite translate,
and $T\in\mathcal C_u$ for one original column.  Its seed is zero
off $N_u(\mathbb A_{F,f})H_u(\mathbb A_{F,f})x_uK$, and on that set is
\begin{equation}\label{global-gram-original-seed}
 \Phi_T(g,nhx_uk)=\psi_{u,f}(n)\chi_2(k)
                       T(W_u(g))(1,h).
\end{equation}
In the next theorem the original decomposition-independence condition
\texttt{PhWellDef} is retained explicitly.  It is exactly a premise
of the source's \texttt{IUnfold} contract.  It concerns the original
finite lattice and phases; replacing them by a convenient level is
not permitted.

\begin{theorem}[The actual weighted folding integral]
\label{global-gram-ifinite}
For every original setting, the full original \texttt{IUnfold}
implication holds.  More explicitly, suppose the prescribed seeds
\eqref{global-gram-original-seed} satisfy \texttt{PhWellDef}, and
the original weighted-seed condition holds.  Then for every original
column $T$ and every compact $\Omega$,
\begin{equation}\label{global-gram-ifinite-eq}
 I_\Omega(T):=
 \int_{S_u(F)\backslash G(\mathbb A_F)}
              m_\Omega(y)|\Phi_T(y)|\,dy<\infty.
\end{equation}
The quotient uses counting measure on $S_u(F)$ and the original
adelic Haar measure; the value is independent of the measurable
fundamental domain.  In particular the original gap hypothesis
supplies the weighted-seed premise through
Theorem~\ref{actual-seed-decay-theorem}.
\end{theorem}

\begin{proof}
The arithmetic quotient $S_u(F)\backslash S_u(\mathbb A_F)$ is
compact.  Indeed, $H_u$ is positive definite at every real place.
Its finite adelic double quotient by a compact open subgroup is
finite: it classifies lattices of one fixed positive definite
Hermitian genus, and finiteness follows from integral positive
definite lattice reduction.  One may use the trace quadratic form;
the matrices giving the action of fixed generators of
$\mathcal O_E$ are bounded on each of the finitely many resulting
quadratic lattices.  A finite level adds only finitely many choices.
Together with compactness of $H_u(F_\infty)$ this proves compactness
of its full arithmetic quotient.  The unipotent group $N_u$ has a
central filtration with additive vector-group quotients, whose
arithmetic quotients are compact.  Induction on that filtration,
then the semidirect product, proves the assertion for $S_u$.
Choose a compact representative set $C_S$.

There is a finite $M_T$ such that
\begin{equation}\label{global-gram-col-bound}
                 |T(\xi)(1,h)|\le M_T\|\xi\|
                 \qquad(h\in H_u(\mathbb A_{F,f})).
\end{equation}
To see this, take the finitely many classes modulo $H_u(F)$ and
the original compact open column level.  Evaluation on each class
is a linear functional on the finite-dimensional $V_{\pi,u}$.
The rational left action is transported through the unitary
representation of $H_u(F_\infty)$, and the right level character
has absolute value one.  The maximum of these finitely many
functional norms gives \eqref{global-gram-col-bound}.
The finite characters in \eqref{global-gram-original-seed} have
absolute value one: their arguments have an integral multiple in
the annihilating integral lattice, so their values are roots of
unity.  Thus $|\Phi_T(g,y_f)|\le M_T\|W_u(g)\|$ on its support.

Let $y=(g,y_f)$ be in that support and fix any $s\in S_u(F_\infty)$.
Write $y_f=s_fx_uk$.  Then
\[
 y=(s^{-1},s_f)(sg,x_uk)
   =\gamma c(1,x_uk)(sg,1)
 \quad(\gamma\in S_u(F),\ c\in C_S).
\]
Return multiplicity is invariant under left multiplication by
$\Gamma$.  Apply Theorem~\ref{global-gram-multiplicity} to the
compact set $C_S(1,x_uK)$ and then take the infimum over $s$.
There is no measurable-choice step and no change of quotient norm:
\begin{equation}\label{global-gram-weight-majorant}
 m_\Omega(y)|\Phi_T(y)|
 \le A M_T\nu_u(g)^{N_X}\|W_u(g)\|
       1_{S_u(\mathbb A_{F,f})x_uK}(y_f).
\end{equation}
The right side is left $S_u(\mathbb A_F)$-invariant.

Use the original $ds_{u,\infty}=d\mu_Nd\mu_H$ and any finite
Haar measure $ds_{u,f}$, and define their quotient measures by
Weil's formula.  The groups in this calculation are unimodular.
For $S_u$ this follows also because the determinant character of
the action of its anisotropic unitary Levi on the unipotent radical
is trivial.  Disintegrating first by $S_u(F)\backslash S_u(\mathbb A_F)$,
then at infinity and at the finite places, bounds
\eqref{global-gram-ifinite-eq} by
\[
 A M_T\operatorname{vol}(S_u(F)\backslash S_u(\mathbb A_F))
       \nu_{u,f}(S_u(\mathbb A_{F,f})x_uK)
       \int_{S_u(F_\infty)\backslash G(F_\infty)}
                \nu_u(g)^{N_X}\|W_u(g)\|\,dq_u(g).
\]
Every factor is finite.  The finite quotient set is the image of
a compact set, and its quotient Haar measure is locally finite.
The product of the two finite-place factors is unchanged if
$ds_{u,f}$ is rescaled.  The fixed infinite-place and adelic
normalizations have not been changed.

Finally the original rational character product and column
equivariance make $\Phi_T$ left $S_u(F)$-invariant; these identities
are applied to the well-defined formula in the hypothesis.
Its measurability follows from the continuous infinite-place seed
and the locally constant finite-level factors.  Therefore Weil's
formula and fundamental-domain independence apply to the actual
seed, proving the original \texttt{IUnfold} contract.
\end{proof}

\begin{corollary}[Absolute Poincar\'e returns and their Gram identity]
\label{global-gram-poincare}
Under the preceding original local decomposition hypothesis, retain
the original gap and finite level.  Set
\[
 P_T(x)=\sum_{\delta\in S_u(F)\backslash\Gamma}\Phi_T(\delta x).
\]
This series converges absolutely at every $x$, independently of its
representatives, and belongs to the original $L^2(\Gamma\backslash G(\mathbb A_F))$.
For two original words and columns, put
$S_0=S_u(F)$, $S_1=S_{u'}(F)$ and
$S_r=S_0\cap rS_1r^{-1}$.  Then
\begin{equation}\label{global-gram-unfolding}
 \int_{\Gamma\backslash G(\mathbb A_F)}P_T(x)\overline{P_{T'}(x)}\,dx
 =\sum_{r\in S_0\backslash\Gamma/S_1}
       \int_{S_r\backslash G(\mathbb A_F)}
          \Phi_T(x)\overline{\Phi_{T'}(r^{-1}x)}\,dx .
\end{equation}
The sum of the corresponding absolute integrals is finite.  Every
double coset and its actual stabilizer occurs; no identity-return
truncation is made.
\end{corollary}

\begin{proof}
Let $\alpha$ be the compact reproducer supplied by
Theorem~\ref{actual-seed-decay-theorem}.  Combining it with
$\overline{\chi_2(k)}1_K(k)/\mu_f(K)$ gives a bounded compact
reproducer for $\Phi_T$ with the actual adelic Haar measure.
The reciprocal level volume is essential.  On a common compact
support $\Omega$ containing $1$, write its absolute bound as $a$.
The positive unfolding of this reproduction inequality gives
\[
 \sum_{\delta\in S_0\backslash\Gamma}|\Phi_T(\delta x)|
 \le a\,m_\Omega(x)\,A_T,
 \qquad A_T=\int_{S_0\backslash G(\mathbb A_F)}|\Phi_T|.
\]
Indeed, after integrating over $S_0\backslash G$ the pairs
$(\delta,s)\mapsto\delta^{-1}s$ run through $\Gamma$ exactly once.
The point count is $m_\Omega(x)$.  Theorem~\ref{global-gram-ifinite}
gives $A_T<\infty$ because $m_\Omega\ge1$, and
Theorem~\ref{global-gram-multiplicity} makes each $m_\Omega(x)$ finite.
Unfolding the positive square majorant gives
$\|P_T\|_2^2\le a A_T I_\Omega(T)<\infty$.
For two seeds the same argument bounds their single unfolded
absolute product by $a A_T I_\Omega(T')$.  Tonelli then permits
the decomposition into all double cosets, with stabilizer $S_r$.
Absolute convergence justifies the complex-valued unfolding and
proves \eqref{global-gram-unfolding} with the displayed conjugation.
\end{proof}

\begin{theorem}[Finite arithmetic fibres of the framed returns]
\label{global-gram-levi-fibres}
Retain one original flag, with isotropic subspace $X$, fixed vector
$z$ modulo $X$, and positive definite complement $W$.  For every
original lattice $\Lambda$ let
\[
 A_\Lambda=\{s\in S_u(F):s\Lambda=\Lambda\},\qquad
 N_\Lambda=A_\Lambda\cap N_u(F).
\]
Then $A_\Lambda/N_\Lambda$ is finite.  More precisely, with $p_W$
the projection associated to the original hyperbolic vectors and $z$,
put
\[
 L_W=p_W(\Lambda\cap X^\perp)\subset W.
\]
The quotient $A_\Lambda/N_\Lambda$ injects into the finite isometry
group of the positive definite rational trace lattice
$(L_W,\operatorname{Tr}_{E/\mathbb Q}h)$.
For each pair of original lattices $\Lambda_1,\Lambda_2$ there is
a finite set $F_{12}\subset S_u(F)$ such that every
$s\in S_u(F)$ with $s\Lambda_2=\Lambda_1$ satisfies
$f^{-1}s\in N_{\Lambda_2}$ for some $f\in F_{12}$.
This is the full original \texttt{LeviFinC} assertion.
\end{theorem}

\begin{proof}
The rational subspace $X^\perp$ is the direct sum
$X\oplus Ez\oplus W$.  The map $p_W$ is rational and fixes $W$.
Consequently $L_W$ is a full $\mathcal O_E$-lattice in $W$:
the intersection $\Lambda\cap X^\perp$ spans $X^\perp$, and its
image is finitely generated and spans $W$.  The denominators of
the projection are retained in $L_W$; no saturation is performed.

An element of $S_u(F)$ acts unipotently on the successive pieces
of the original isotropic flag, fixes the class of $z$ in
$X^\perp/X$, and induces a unitary transformation $s_W$ on $W$.
For $x\in X^\perp$ one has $p_W(sx)=s_Wp_W(x)$.  Thus
$s\Lambda=\Lambda$ implies $s_WL_W=L_W$.  Conversely the kernel
of $s\mapsto s_W$ on $S_u(F)$ is exactly its prescribed unipotent
radical: it acts trivially on $X^\perp/X$, including its $Ez$
summand, as well as on every graded piece of the flag.  Intersecting
with $A_\Lambda$ gives precisely the original subgroup
$N_\Lambda$, with the original lattice condition intact.

Choose a $\mathbb Z$-basis $b_1,\ldots,b_d$ of $L_W$.  After
clearing one fixed denominator, its trace form is integral and
positive definite.  An isometry sends $b_i$ to a lattice vector
of the fixed norm of $b_i$.  Each such norm sphere meets the lattice
in finitely many points, so only finitely many tuples of images
are possible.  This proves the asserted finiteness without an
estimate uniform in $d$, the discriminant or the projected lattice.

If no $s$ sends $\Lambda_2$ to $\Lambda_1$, take $F_{12}$ empty.
Otherwise fix $s_0$ doing so, choose representatives $T$ for the
finite quotient $A_{\Lambda_2}/N_{\Lambda_2}$, and set
$F_{12}=s_0T$.  Since $s_0^{-1}s\in A_{\Lambda_2}$,
the asserted right-coset factorization follows.  The description
of $N_\Lambda$ above is exactly the flag and orthogonal-quotient
definition of \texttt{NLam}, so it proves the stated source
contract, not a quotient by a larger unipotent group.
\end{proof}

This closes the actual arithmetic \texttt{MultPoly},
\texttt{IUnfold} and \texttt{LeviFinC} suppliers, and proves the
displayed Poincar\'e identity under the explicit original local
decomposition hypothesis.
It does not yet prove the finer original \texttt{OOLintSumC}
statement for all framed returns, the class-frame identification
\texttt{SigRowsC}, or the normalized \texttt{RowFormC} formula.
Their finite level weights, phases, torsors and Haar constants still
have to be matched.  No bound uniform in $X$ follows from the
finite constants above.  In particular this is not yet the required
positive-trace estimate or an aggregate nonvanishing theorem.
